\section{Research Question and Hypothesis}

\subsection{Research Question}

The central research question is:

\begin{quote}
To what extent can a lightweight two-dimensional axisymmetric Python solver, coupling sparse Laplace/Poisson electrostatics with electrostatic--capillary residual analysis and effective space-charge shielding closures, predict Taylor-cone geometry and onset-voltage trends on consumer-grade hardware?
\end{quote}

\subsection{Hypothesis}

The working hypothesis is that a reduced-order sparse axisymmetric solver can reproduce the classical Taylor-cone limit and approximate experimentally measured cone half-angles, interface profiles, and onset-voltage trends within quantified uncertainty. The expected dominant error sources are finite-grid boundary representation, simplified effective space-charge closures, and physical effects not included in the static model, such as liquid flow, current transport, and full leaky-dielectric dynamics.

\subsection{Scope}

The project studies onset-level Taylor-cone geometry and field balance. It does not aim to simulate the full downstream cone-jet, droplet breakup, emission-current law, plume dynamics, or complete thruster performance. This limited scope is intentional: the purpose is to test whether a transparent lightweight model captures key observables before using more complex simulations.
